\section{Light --- Compiler --- Start Up}
%%% generated by the blueprint pipeline from TCSlib.Programs.Light.Compiler.StartUp
\subsection*{Light --- Compiler --- Start Up: Introduction}

Declarations of the Lean module \texttt{TCSlib.Programs.Light.Compiler.StartUp}.

\subsection*{Light --- Compiler --- Start Up: Results}

\begin{theorem}[Effects of head initialization code]
\lean{Light.Compiler.effects_headStraight}
\label{Light.Compiler.effects_headStraight}
\leanok
\uses{Light.Compiler.cNEG, Light.Compiler.cONE, Light.Compiler.cZERO, Light.Compiler.headStraight, WordRAM.effect_one, WordRAM.effect_sub, WordRAM.effects, WordRAM.effects_cons, WordRAM.effects_nil, WordRAM.wd}
\difficulty{3}
For every initial memory, executing the head-initialization code yields the memory
obtained by setting the designated constant-one cell to the word value $1$, the zero
cell to $0$, and the negative-one cell to $-1$, in that order.
\end{theorem}

\begin{theorem}[Initialized constant cells]
\lean{Light.Compiler.effects_headStraight_consts}
\label{Light.Compiler.effects_headStraight_consts}
\leanok
\uses{Light.Compiler.cNEG, Light.Compiler.cONE, Light.Compiler.cZERO, Light.Compiler.effects_headStraight, Light.Compiler.headStraight, WordRAM.effects, WordRAM.wd}
\difficulty{3}
For any word size $W$ and initial memory function $m : \mathbb{Z} \to \{0,1\}^W$,
executing the three-instruction initialization sequence leaves the cells at addresses
$-5$, $-7$, and $-9$ containing the word representations of $1$, $0$, and $-1$,
respectively.
\end{theorem}

\begin{theorem}[Successor formula for pool code]
\lean{Light.Compiler.poolCode_succ}
\label{Light.Compiler.poolCode_succ}
\leanok
\uses{Light.Compiler.cONE, Light.Compiler.cPool, Light.Compiler.poolCode}
\difficulty{1}
For every natural number $N$, the pool initialization code for $N+1$ entries is the pool
initialization code for $N$ entries followed by the instruction that adds the constant
one to the value at pool address $c_{\mathrm{Pool}}(N)$ and stores the result at
$c_{\mathrm{Pool}}(N+1)$, where $c_{\mathrm{Pool}}(n)=-(25+4n)$.
\end{theorem}

\begin{theorem}[Memory preserved outside the pool]
\lean{Light.Compiler.agreeOutside_poolCode}
\label{Light.Compiler.agreeOutside_poolCode}
\leanok
\uses{Light.Compiler.Pool, Light.Compiler.cPool, Light.Compiler.poolCode, WordRAM.AgreeOutside, WordRAM.agreeOutside_effects, WordRAM.effects}
\difficulty{4}
Let $W$ be a word width, $N\in\bbn$, and let $m:\bbr\to\text{BitVec}(W)$ be a memory.
The pool-initialization program sets the pool cell at index $0$ and then fills the cells
at indices $1$ through $N$; after executing it, the resulting memory agrees with $m$ at
every integer address outside the set of pool addresses $-(25+4n)$ for $0\le n\le N$.
\end{theorem}

\begin{theorem}[Pool initialization]
\lean{Light.Compiler.effects_poolCode}
\label{Light.Compiler.effects_poolCode}
\leanok
\uses{Light.Compiler.agreeOutside_poolCode, Light.Compiler.cONE, Light.Compiler.cPool, Light.Compiler.poolCode, Light.Compiler.poolCode_succ, WordRAM.AgreeOutside.read, WordRAM.effect_add, WordRAM.effect_sub, WordRAM.effects, WordRAM.effects_append, WordRAM.effects_cons, WordRAM.effects_nil, WordRAM.wd}
\difficulty{5}
Let $W,N\in\mathbb{N}$, let $m:\mathbb{Z}\to\mathrm{BitVec}(W)$ be a memory state with
$m(c_{\mathrm{one}})=1$ as a $W$-bit word, and let $c_{\mathrm{pool}}(n)=-(25+4n)$ and
$c_{\mathrm{one}}=-5$. The instruction sequence that first sets the pool cell at index
$0$ to $1$ and then successively adds $1$ to each preceding pool cell leaves every cell
with index $n\le N$ equal to the $W$-bit representation of $n$.
\end{theorem}

\begin{definition}[Pool memory transformation]
\lean{Light.Compiler.poolMem}
\label{Light.Compiler.poolMem}
\leanok
\uses{Light.Compiler.headStraight, Light.Compiler.poolCode, WordRAM.effects}
For a natural number $N$ and an integer-addressed memory $m$ whose values are $W$-bit
words, define the pool memory to be the result of applying the effects of the pool code
for $N$ to the memory obtained by applying the effects of the head-straight code to $m$.
\end{definition}

\begin{theorem}[Memory preserved outside the constants]
\lean{Light.Compiler.agreeOutside_poolMem}
\label{Light.Compiler.agreeOutside_poolMem}
\leanok
\uses{Light.Compiler.Constants, Light.Compiler.agreeOutside_poolCode, Light.Compiler.effects_headStraight, Light.Compiler.poolMem, WordRAM.AgreeOutside, WordRAM.AgreeOutside.refl, WordRAM.AgreeOutside.trans_union, WordRAM.AgreeOutside.update}
\difficulty{3}
Let $W$ be a word width, $N\in\bbn$, and let $m:\bbr\to\text{BitVec}(W)$ be a memory.
First initialize the cells at addresses $-5$, $-7$, and $-9$ to $1$, $0$, and $-1$,
respectively, and then initialize the pool cells at addresses $-(25+4n)$ for $0\le n\le
N$. The resulting memory agrees with $m$ at every integer address outside the union of
those three constant addresses and the pool addresses.
\end{theorem}

\begin{theorem}[Initialized constant registers]
\lean{Light.Compiler.poolMem_consts}
\label{Light.Compiler.poolMem_consts}
\leanok
\uses{Light.Compiler.agreeOutside_poolCode, Light.Compiler.cNEG, Light.Compiler.cONE, Light.Compiler.cZERO, Light.Compiler.effects_headStraight_consts, Light.Compiler.headStraight, Light.Compiler.poolMem, WordRAM.AgreeOutside.read, WordRAM.effects, WordRAM.wd}
\difficulty{3}
For every natural number $N$ and every machine memory $m$ with $W$-bit values, after the
straight-line initialization of the constant registers and pool, the register at the one
address contains the $W$-bit representation of $1$, the register at the zero address
contains the representation of $0$, and the register at the negative-one address
contains the representation of $-1$.
\end{theorem}

\begin{theorem}[Initialized pool values]
\lean{Light.Compiler.poolMem_pool}
\label{Light.Compiler.poolMem_pool}
\leanok
\uses{Light.Compiler.cPool, Light.Compiler.effects_headStraight_consts, Light.Compiler.effects_poolCode, Light.Compiler.poolMem, WordRAM.wd}
\difficulty{1}
For every natural number $N$, every machine memory $m$ with $W$-bit values, and every
natural number $n\le N$, after the straight-line initialization of the constant
registers and pool, the value at pool address $c_{\mathrm{Pool}}(n)=-(25+4n)$ is the
$W$-bit representation of $n$.
\end{theorem}

\begin{theorem}[Successive zero-code blocks]
\lean{Light.Compiler.zeroCode_succ}
\label{Light.Compiler.zeroCode_succ}
\leanok
\uses{Light.Compiler.cONE, Light.Compiler.cStack, Light.Compiler.mainFrame, Light.Compiler.zeroCode}
\difficulty{1}
For every natural number $F$, the code that initializes $F+1$ stack cells is the code
that initializes the first $F$ cells followed by the instruction that subtracts the
value in the constant-one cell from itself and stores the result in the stack cell
indexed by the main-frame offset plus $F$.
\end{theorem}

\begin{theorem}[Writes of the zeroing code]
\lean{Light.Compiler.writesIn_zeroCode}
\label{Light.Compiler.writesIn_zeroCode}
\leanok
\uses{Light.Compiler.MainFrame, Light.Compiler.cStack, Light.Compiler.mainFrame, Light.Compiler.zeroCode, WordRAM.WritesIn}
\difficulty{3}
For every natural number $F$, each instruction in the zeroing code for $F$ writes only
to an address in the main frame of size $F$.
\end{theorem}

\begin{theorem}[Zeroing the Stack Cells]
\lean{Light.Compiler.effects_zeroCode}
\label{Light.Compiler.effects_zeroCode}
\leanok
\uses{Light.Compiler.cONE, Light.Compiler.cStack, Light.Compiler.mainFrame, Light.Compiler.writesIn_zeroCode, Light.Compiler.zeroCode, Light.Compiler.zeroCode_succ, WordRAM.AgreeOutside.read, WordRAM.agreeOutside_effects, WordRAM.effect_sub, WordRAM.effects, WordRAM.effects_append, WordRAM.effects_cons, WordRAM.effects_nil, WordRAM.wd}
\difficulty{5}
Let $W$ and $F$ be natural numbers, and let $m$ be a memory state assigning a $W$-bit
word to each integer address. If the word at address $-5$ is the $W$-bit representation
of $1$, then for every natural number $x<F$, executing the program that subtracts the
word at address $-5$ from itself and stores the result at address $-2(4+x)$, once for
each index from $0$ through $F-1$, leaves the word at address $-2(4+x)$ equal to the
$W$-bit representation of $0$.
\end{theorem}

\begin{theorem}[Memory preserved outside the main frame]
\lean{Light.Compiler.agreeOutside_frameCode}
\label{Light.Compiler.agreeOutside_frameCode}
\leanok
\uses{Light.Compiler.MainFrame, Light.Compiler.cFP, Light.Compiler.cStack, Light.Compiler.frameCode, Light.Compiler.mainFrame, Light.Compiler.writesIn_zeroCode, WordRAM.AgreeOutside, WordRAM.WritesIn.mono, WordRAM.agreeOutside_effects, WordRAM.effects}
\difficulty{4}
Let $W$ be a word width, $m:\mathbb{Z}\to\mathrm{BitVec}(W)$ a memory, and $F$ a
natural number with $3\le F$. Executing the outermost-frame initialization code changes
memory only at the frame-pointer cell $c_{\mathrm{FP}}=-11$ and at the frame cells
$-2j$ for $4\le j<4+F$; equivalently, the initial and resulting memories agree outside
${c_{\mathrm{FP}}}\cup\{-2j:4\le j<4+F\}$.
\end{theorem}

\begin{theorem}[Effects of frame code]
\lean{Light.Compiler.effects_frameCode}
\label{Light.Compiler.effects_frameCode}
\leanok
\uses{Light.Compiler.MainFrame, Light.Compiler.cARG1, Light.Compiler.cARG2, Light.Compiler.cFP, Light.Compiler.cPool, Light.Compiler.cStack, Light.Compiler.cZERO, Light.Compiler.frameCode, Light.Compiler.mainFrame, Light.Compiler.writesIn_zeroCode, Light.Compiler.zeroCode, WordRAM.AgreeOutside.cell, WordRAM.agreeOutside_effects, WordRAM.effect_add, WordRAM.effect_sub, WordRAM.effects, WordRAM.effects_append, WordRAM.effects_cons, WordRAM.effects_nil, WordRAM.wd}
\difficulty{5}
Suppose the initial memory holds $0$ at the zero address, the value $2\cdot 4$ at the
pool cell for offset $2\cdot 4$, and values $i$ and $j$ at the two argument addresses.
For every natural number $F$, executing the frame code produces the memory obtained by
first setting the frame pointer to the stack address for index $4$, then executing the
code that zeros the first $F$ local stack cells, and finally setting the next two stack
cells to $i$ and $j$, in that order.
\end{theorem}

\begin{definition}[Start memory transformation]
\lean{Light.Compiler.startMem}
\label{Light.Compiler.startMem}
\leanok
\uses{Light.Compiler.headStraight, Light.Compiler.startStraight, WordRAM.effects}
For natural numbers $F$ and $N$ and an integer-addressed memory $m$ whose values are
$W$-bit words, define the start memory to be the result of applying the effects of the
start-straight code for $F$ and $N$ to the memory obtained by applying the effects of
the head-straight code to $m$.
\end{definition}

\begin{theorem}[Starting memory decomposition]
\lean{Light.Compiler.startMem_eq}
\label{Light.Compiler.startMem_eq}
\leanok
\uses{Light.Compiler.frameCode, Light.Compiler.poolMem, Light.Compiler.startMem, WordRAM.effects, WordRAM.effects_append}
\difficulty{1}
For all natural numbers $F,N$ and every WordRAM memory $m$, initializing the compiler's
starting memory with parameters $F,N$ gives the same memory as first applying the head
and pool initialization to $m$, then applying the frame initialization with parameter
$F$.
\end{theorem}

\begin{theorem}[Memory preserved outside startup locations]
\lean{Light.Compiler.agreeOutside_startMem}
\label{Light.Compiler.agreeOutside_startMem}
\leanok
\uses{Light.Compiler.Constants, Light.Compiler.MainFrame, Light.Compiler.agreeOutside_frameCode, Light.Compiler.agreeOutside_poolMem, Light.Compiler.cFP, Light.Compiler.startMem, Light.Compiler.startMem_eq, WordRAM.AgreeOutside, WordRAM.AgreeOutside.trans_union}
\difficulty{3}
Let $W$ be a word width, $F,N\in\bbn$ with $3\le F$, and let $m:\bbr\to\text{BitVec}(W)$
be a memory. The startup program initializes the constants at $-5,-7,-9$, the pool cells
at addresses $-(25+4n)$ for $0\le n\le N$, the frame pointer at $-11$, and the frame
cells at addresses $-2j$ for $4\le j<4+F$, also copying the two argument values into
frame cells. After executing it, the resulting memory agrees with $m$ at every integer
address outside the union of the constant and pool addresses, the frame-pointer address,
and the frame-cell addresses.
\end{theorem}

\begin{theorem}[Relation after initialization]
\lean{Light.Compiler.rel_startMem}
\label{Light.Compiler.rel_startMem}
\leanok
\uses{Light.Compiler.Rel, Light.Compiler.Sizes, Light.Compiler.agreeOutside_frameCode, Light.Compiler.agreeOutside_poolMem, Light.Compiler.cARG1, Light.Compiler.cARG2, Light.Compiler.cStack, Light.Compiler.effects_frameCode, Light.Compiler.effects_zeroCode, Light.Compiler.frameCode, Light.Compiler.mainFrame, Light.Compiler.poolMem, Light.Compiler.poolMem_consts, Light.Compiler.poolMem_pool, Light.Compiler.startMem, Light.Compiler.startMem_eq, Light.Compiler.writesIn_zeroCode, Light.frame, WordRAM.AgreeOutside.cell, WordRAM.AgreeOutside.read, WordRAM.agreeOutside_effects, WordRAM.effects, WordRAM.wd}
\difficulty{5}
Let $W$ be a natural number and let a compiler size configuration have frame width $F$,
dispatch position $D$, and space limit $S$. Suppose $3\le F$ and $8\le D$. Let $\mu$
assign an integer to each natural-number address, and let $i,j\in\mathbb Z$. If the
initial machine memory at every address $a<S$ is the $W$-bit encoding of $\mu(a)$, and
the two argument registers contain the $W$-bit encodings of $i$ and $j$, respectively,
then after the initialization program is run, the machine is related at stack index $4$
to the light state with memory $\mu$ and local-variable values $0,i,j$ at indices
$0,1,2$, and zero at every other index. Here the compiler size configuration includes
the requirement that its word size fits its stated limits, pool, and stack.
\end{theorem}

\begin{theorem}[Straight instructions in the initialization prefix]
\lean{Light.Compiler.straight_headStraight}
\label{Light.Compiler.straight_headStraight}
\leanok
\uses{Light.Compiler.headStraight, WordRAM.Straight}
\difficulty{1}
Each instruction in the fixed initialization prefix is straight.
\end{theorem}

\begin{theorem}[Straightness of the initial code]
\lean{Light.Compiler.straight_startStraight}
\label{Light.Compiler.straight_startStraight}
\leanok
\uses{Light.Compiler.frameCode, Light.Compiler.poolCode, Light.Compiler.startStraight, Light.Compiler.zeroCode, WordRAM.Straight}
\difficulty{1}
For all natural numbers $F$ and $N$, every instruction in the initial code formed from
the pool initialization and frame initialization is a constant assignment, addition,
subtraction, multiplication, load, or store instruction.
\end{theorem}

\begin{theorem}[Straight start code length]
\lean{Light.Compiler.length_startStraight}
\label{Light.Compiler.length_startStraight}
\leanok
\uses{Light.Compiler.frameCode, Light.Compiler.poolCode, Light.Compiler.startStraight, Light.Compiler.zeroCode}
\difficulty{2}
For all natural numbers $F$ and $N$, the straight start code formed from the pool code
and frame code has length $N+F+4$.
\end{theorem}

\begin{theorem}[Startup steps to the main procedure]
\lean{Light.Compiler.steps_start}
\label{Light.Compiler.steps_start}
\leanok
\uses{Light.Compiler.agreeOutside_frameCode, Light.Compiler.cNEG, Light.Compiler.codeAt_headCode, Light.Compiler.codeAt_startCode, Light.Compiler.dispPos, Light.Compiler.effects_headStraight_consts, Light.Compiler.frameSize, Light.Compiler.headStraight, Light.Compiler.length_startStraight, Light.Compiler.mainPos, Light.Compiler.poolMem_consts, Light.Compiler.startMem, Light.Compiler.startMem_eq, Light.Compiler.steps_jump, Light.Compiler.straight_headStraight, Light.Compiler.straight_startStraight, Light.Compiler.three_le_frameSize, Light.Program, Light.compileProgram, WordRAM.AgreeOutside.cell, WordRAM.CodeAt.left, WordRAM.CodeAt.right, WordRAM.InRange, WordRAM.Steps, WordRAM.Steps.cast_count, WordRAM.Steps.trans, WordRAM.steps_straight, WordRAM.wd}
\difficulty{4}
Let $P$ be a light program, let $p_0$ be a procedure index, let $d$ be a Boolean flag,
and let $W$ be a word size. Suppose the integer $-1$ satisfies $-2^W \le 2(-1) < 2^W$,
and let $m$ be any initial word-RAM memory. Executing the compiled program for
$\mathrm{dispPos}(P,d)+\mathrm{frameSize}(P)+9$ steps from program counter $0$ and
memory $m$ reaches program counter $\mathrm{mainPos}$ with memory given by applying the
startup initialization to $m$.
\end{theorem}
